% Chapter 10 --- Beyond the paper: a call-by-name PPL on SCones, the
% diagonal bilinear stability bridge, and the CBN headline
% ex_geom_CBN_mass_one.
%
% Mirrors:
%   theories/programs/ppl_cbn.v        the CBN trunk: tyD_CBN / ctxD_CBN
%                                      pure-fragment eD_CBN (var/tt/pair/
%                                      fst/snd/lam/app/let), ne_fix via the
%                                      §9.2 SCones operator Yfix
%   theories/programs/ppl_cbn_eff.v    CBN effects: sample/score/real/add/
%                                      mul clauses (option (γ): degenerate
%                                      add/mul as the headlines baseline)
%   theories/programs/ppl_cbn_bool.v   CBN boolean cascade: true/false/
%                                      Bernoulli/if via sc_to_sh and the
%                                      bool_case_scones promotion
%   theories/programs/ppl_cbn_arith.v  CBN add_FMeas / mul_FMeas math
%                                      foundation (the genuine arithmetic
%                                      on the FMeas-cone)
%   theories/programs/ppl_cbn_headlines.v
%                                      structural _denot_E reductions for
%                                      the CBN headlines under option (γ)
%   theories/programs/ppl_cbn_geom.v   THE CBN HEADLINE
%                                      ex_geom_CBN_mass_one : the CBN
%                                      denotation of the geometric program
%                                      has total mass 1 (axiom-free).  Built
%                                      via the SCones endomorphism
%                                      phi_CBN_geom and the Kleene cascade
%                                      mass identity 1 - (1/2)^n
%   theories/stable/diag_bilinear_tensor.v
%                                      M1 linhom_to_stablehom + meas-stable,
%                                      M2 meas_stable_comp_post/pre,
%                                      M3 id_spair_meas_stable /
%                                          spair_meas_stable,
%                                      M4 meas_stable_diag_bilinear_tensor
%                                      (the headline: structural composition
%                                      via Thm 5.12 tensor-curry).
%
% Status note: the WHOLE chapter is AXIOM-FREE — only the three classical
% boolp axioms.  The two centrepieces are
%   - ex_geom_CBN_mass_one (ppl_cbn_geom.v): the CBN-side parallel of the
%     CBV ex_geom_cbv_mass_one of Chapter 9 (ex_reject_headline.v),
%     closed via the unary
%     "shift-by-Dirac" insight (the ELSE branch [|1|] + #"g" @ () is the
%     unary linear pushforward through (+ 1), an icones_hom, no bilinear
%     bridge needed at the headline);
%   - meas_stable_diag_bilinear_tensor (diag_bilinear_tensor.v): the
%     SCones↔ICones-tensor bilinear stability bridge, structural unblocker
%     of generic CBV recursion at [eD ne_fix], honest CBN bilinear
%     arithmetic, and the CBV/CBN soundness comparison.

%%%%%%%%%%%%%%%%%%%%%%%%%%%%%%%%%%%%%%%%%%%%%%%%%%%%%%%%%%%%%%%%%%%%%%%%
\chapter{Beyond the paper: a call-by-name model, the diagonal bilinear
stability bridge, and the CBN headline}\label{ch:cbn}

\paragraph{This chapter is beyond the Ehrhard--Geoffroy paper.}
The paper stops at the \S9 LNL/Seely model. The previous chapter built a
\emph{call-by-value} interpretation of a higher-order PPL on the
Eilenberg--Moore category $\mathrm{EM}(\mathord{!})$. This chapter builds
the parallel \emph{call-by-name} interpretation of the \emph{same}
surface language on the cartesian closed $\mathbf{SCones}$
(\autoref{thm:scones-ccc}) --- and the structural unblocker that ties
the two readings together at the level of bilinear-into-tensor
arithmetic, the diagonal bilinear stability bridge
\rocq{Icones.stable.diag_bilinear_tensor.meas_stable_diag_bilinear_tensor}.
The chapter closes with the axiom-free CBN headline
\rocq{Icones.programs.ppl_cbn_geom.ex_geom_CBN_mass_one}: the call-by-name
denotation of the geometric program has total mass $1$, the
centrepiece of this chapter. (The CBV-side twin now exists:
\rocq{Icones.programs.ex_reject_headline.ex_geom_cbv_mass_one} of
\autoref{thm:cbv-ex-geom}, proved through the seeded value-fixpoint
combinator of \autoref{sec:cbv-yfix-fun}, so both interpretations of
the shared surface program prove mass one for the geometric counter.)

\paragraph{Where the material lives.}
\texttt{theories/programs/ppl\_cbn.v} (the CBN trunk: type / context
translations $\sem{-}_n$ and the pure fragment of $\mathtt{eD\_CBN}$),
\texttt{theories/programs/ppl\_cbn\_eff.v} (the CBN effect clauses for
$\mathtt{ne\_sample}/\mathtt{ne\_score}/\mathtt{ne\_real}/\mathtt{ne\_add}
/\mathtt{ne\_mul}$ under the degenerate-arithmetic option $\gamma$),
\texttt{theories/programs/ppl\_cbn\_bool.v} (the CBN boolean cascade
via the $\texttt{sc\_to\_sh}$ promotion of paper Lemma 7.31),
\texttt{theories/programs/ppl\_cbn\_arith.v} (the genuine arithmetic
$\mathtt{add\_FMeas}/\mathtt{mul\_FMeas}$ math foundation on
$\FMeas$), \texttt{theories/programs/ppl\_cbn\_headlines.v} (the
structural $\_\mathtt{denot}\_E$ reductions of the QBS-style headlines
under option $\gamma$),
\texttt{theories/programs/ppl\_cbn\_geom.v} (the centrepiece headline
\rocq{Icones.programs.ppl_cbn_geom.ex_geom_CBN_mass_one}), and
\texttt{theories/stable/diag\_bilinear\_tensor.v} (the second centrepiece,
the diagonal bilinear stability bridge
\rocq{Icones.stable.diag_bilinear_tensor.meas_stable_diag_bilinear_tensor}).

%%%%%%%%%%%%%%%%%%%%%%%%%%%%%%%%%%%%%%%%%%%%%%%%%%%%%%%%%%%%%%%%%%%%%%%%
\section{The call-by-name interpretation}\label{sec:cbn-interp}

The CBN interpretation reuses the entire surface syntax
\rocq{Icones.programs.ppl.named_expr} of \autoref{sec:cbv-ppl}: the same
intrinsically-typed multi-variable named-variable inductive denotes both
in CBV (via $\mathrm{EM}(\mathord{!})$) and in CBN (via
$\mathbf{SCones}$). The reading is a step-for-step transport of
Melli\`es' co-Kleisli rule for $\mathord{!}$: types become $\ICones$
objects \emph{without} the $\mathord{!}$-coalgebra wrap, contexts are
$\mathbf{SCones}$ products, every term denotes a stable measurable
function out of its context, and recursion at function type is the
paper-\S9.2 fixpoint operator $\mathrm{Yfix}$ applied at the source-language
function type.

\begin{definition}[The CBN type and context translations]\label{def:cbn-tyD}
\rocq{Icones.programs.ppl_cbn.tyD_CBN}\rocqok
\rocq{Icones.programs.ppl_cbn.ctxD_CBN}\rocqok
\uses{def:scones-hom, thm:scones-prod, thm:stablehom-icone,
def:cbv-bool-cone}
The CBN type translation $\sem{-}_n : \mathtt{ppl\_type} \to \ICones$ is
\[
  \sem{\mathtt{tunit}}_n = \top,\quad
  \sem{\mathtt{tbase}\,X}_n = \FMeas(X),\quad
  \sem{\mathtt{tprod}\,s_1\,s_2}_n = \sem{s_1}_n \mathbin{\&} \sem{s_2}_n,
\]
\[
  \sem{\mathtt{tfun}\,A\,B}_n
    = \sem{A}_n \Rightarrow_s \sem{B}_n,\qquad
  \sem{\mathtt{tbool}}_n = \mathtt{bool\_cone\_car},
\]
and the context translation $\sem{\Gamma}_n$ is the right-associated
$\mathbf{SCones}$ product $\sem{t_1}_n \mathbin{\&} \cdots \mathbin{\&}
\sem{t_n}_n \mathbin{\&} \top$. There is \emph{no} $T = \widetilde{!}
\circ U$ wrap on the codomain of $\sem{\mathtt{tfun}\,A\,B}_n$: in CBN,
effects live inside the type interpretation through the cone of
sub-probability measures $\FMeas$, in the manner of the lazy
QBS reading.
\end{definition}

\begin{remark}[The pragmatic QBS reading on base types]
The CBN translation $\sem{\mathtt{tbase}\,X}_n = \FMeas(X)$ is the
``pragmatic QBS reading'': base values \emph{are} measures, with the
constant Dirac path realising the inclusion $X \hookrightarrow \FMeas(X)$.
This deliberately differs from the Bang-promoted alternative
$\sem{\mathtt{tbase}\,X}_n = \mathord{!}\FMeas(X)$ which the strict
LNL co-Kleisli rule prescribes: it lets the CBN-side $\FMeas$ arithmetic
of \autoref{sec:cbn-arith} fire as a plain $\mathbf{SCones}$ endomorphism
of $\FMeas$, without crossing the $\mathord{!}$/$\mathbf{SCones}$
boundary, and is what makes the CBN headline of
\autoref{thm:cbn-ex-geom} closable in 30 lines.
\end{remark}

\begin{definition}[The CBN term interpretation $\mathtt{eD\_CBN}$]\label{def:cbn-eD}
\rocq{Icones.programs.ppl_cbn.eD_CBN}\rocqok
\rocq{Icones.programs.ppl_cbn.eD_CBN_fix_E}\rocqok
\uses{def:cbn-tyD, thm:scones-ccc, thm:scones-prod}
The CBN denotation
$\mathtt{eD\_CBN} : \mathtt{named\_expr}\,\Gamma\,t \to
\mathbf{SCones}(\sem{\Gamma}_n,\,\sem{t}_n)$ is defined by
structural recursion on $\mathtt{named\_expr}$ (\autoref{sec:cbv-ppl}):
$\mathtt{ne\_var}$ is the projection
\rocq{Icones.programs.ppl_cbn.var_lookup_CBN}; pairing /
projections are $\mathbf{SCones}$ \texttt{spair} / \texttt{sfst} /
\texttt{ssnd}; $\mathtt{ne\_lam}$ is the cartesian-closed $\texttt{curry}$;
$\mathtt{ne\_app}$ is $\mathrm{Ev}$ after $\texttt{spair}$;
$\mathtt{ne\_let}\,x\,M\,K$ is plain $\mathbf{SCones}$ composition
$\sem{K}_n \circ \texttt{spair}\,\mathrm{id}\,\sem{M}_n$
(no Kleisli, no monad: the captured value is passed by the cartesian
pair through the extended product context); the recursion
$\mathtt{ne\_fix}$ unfolds to
\[
  \mathtt{eD\_CBN}(\mathtt{ne\_fix}\,s\,M)
  = \mathrm{Yfix}_{\sem{\mathtt{tfun}\,t_1\,t_2}_n} \circ
     \mathtt{curry}\,\sem{M}_n
\]
(\rocq{Icones.programs.ppl_cbn.eD_CBN_fix_E}), the paper-\S9.2 SCones
fixpoint operator \rocq{Icones.stable.fixpoint.Yfix} applied at the
function-type internal hom $\sem{\mathtt{tfun}\,t_1\,t_2}_n$. The
effect clauses for $\mathtt{ne\_sample}/\mathtt{ne\_score}/\mathtt{ne\_real}
/\mathtt{ne\_add}/\mathtt{ne\_mul}$ are pluggable: the trunk
\texttt{ppl\_cbn.v} abstracts them as section hypotheses, the M3 file
\texttt{ppl\_cbn\_eff.v} discharges them under option $\gamma$ (degenerate
arithmetic, kept lightweight as the headlines baseline), and the
genuine $\FMeas$ arithmetic of \autoref{sec:cbn-arith} is the refinement
direction towards which the diagonal bilinear stability bridge points.
\end{definition}

\begin{lemma}[Free recursion at every function type]\label{lem:cbn-recursion}
\rocq{Icones.programs.ppl_cbn.eD_CBN_fix_E}\rocqok
\uses{def:cbn-eD}
At every function type $\mathtt{tfun}\,A\,B$, $\mathtt{ne\_fix}$ admits a
sound interpretation as a single line of paper-\S9.2 content, with the
fixpoint equation closing in four lines via
\rocq{Icones.stable.fixpoint.Yfix_fix}: this is available
\emph{unconditionally} at any function type $\sem{\mathtt{tfun}\,A\,B}_n =
\sem{A}_n \Rightarrow_s \sem{B}_n$, no bilinear-bridge obligation
arising. (The CBV-side value-fixpoint analogue is currently a
\texttt{precone\_zero} placeholder in \texttt{ppl\_cbv.v} pending a
Kleene port to the clean \texttt{linhom} cone.)
\end{lemma}

\begin{remark}[The CBN boolean cascade]\label{rem:cbn-bool}
\rocq{Icones.programs.infra.bool_case_scones.bool_case_scones}\rocqok
\rocq{Icones.programs.ppl_cbn_bool.cbn_if_clause_def}\rocqok
The boolean cascade in CBN reuses the 2-point sub-probability cone
\rocq{Icones.programs.infra.bool_cone.bool_cone_car} of
\autoref{def:cbv-bool-cone} and its co-pairing
\rocq{Icones.programs.infra.bool_cone.bool_case}, promoted to
$\mathbf{SCones}$ via the dereliction $\mathrm{Der}$ of
\autoref{lem:ders}. The CBN $\mathtt{if}$ clause
\rocq{Icones.programs.ppl_cbn_bool.cbn_if_clause_def} closes by the
same $\texttt{spair}$-then-evaluation trick as
\autoref{def:cbn-eD}, on the 2-point co-pairing realised as an
$\mathbf{SCones}$ morphism through
\rocq{Icones.programs.infra.bool_case_scones.bool_case_scones}.
\end{remark}

%%%%%%%%%%%%%%%%%%%%%%%%%%%%%%%%%%%%%%%%%%%%%%%%%%%%%%%%%%%%%%%%%%%%%%%%
\section{The diagonal bilinear stability bridge}\label{sec:diag-bilin}

\paragraph{Why the bridge.}
A bilinear-into-tensor map $\Phi : G \otimes A \to B$ in $\ICones$ is
not the same data as a stable map: it has type
$\icones_{\mathrm{hom}}\,(G \otimes A)\,B$, in particular it is
\emph{linear} in its first slot. The diagonal evaluation
$g \mapsto \Phi(g \otimes K(g))$ for a measurable stable
$K : G \to A$ is therefore \emph{not} immediately stable, because the
joint $g \mapsto (g, K(g))$ is non-linear in $g$ and Lemma 7.27
(\autoref{lem:lin-totmono}) wants linearity in the first slot of $\Phi$
\emph{viewed as a map out of the cartesian product}, not the tensor.
The recipe of paper Lemma 9.4 (\autoref{lem:stab-lin-swap}) crossed with
the SAFT-discharged tensor--hom iso of \autoref{thm:thm512} repairs the
mismatch by transporting the bilinear input to a curried form, on which
elementwise reasoning works. This is the content of the bridge.

\begin{theorem}[Diagonal bilinear stability bridge]\label{thm:diag-bilin}
\rocq{Icones.stable.diag_bilinear_tensor.meas_stable_diag_bilinear_tensor}\rocqok
\uses{thm:thm512, lem:linhom-to-stablehom, lem:meas-stable-comp,
lem:spair-meas-stable, lem:ev}
For integrable cones $G, E_A, E_B \in \ICones$, a measurable stable map
$K : G \to E_A$, and a bilinear $\Phi :
\icones_{\mathrm{hom}}\,(G \otimes E_A)\,E_B$, the diagonal evaluation
\[
  g \;\mapsto\; \Phi(g \otimes_p K(g))
  \;:\; G \to E_B
\]
is measurable stable. In particular, with $E_A = \mathord{!}A$ and
$E_B = \mathord{!}B$, the recipe applies at the Bang-level CBV
$\mathtt{ne\_fix}$ generic argument and at the
diagonal-into-tensor arithmetic of refined CBN $\mathtt{add}/\mathtt{mul}$.
\end{theorem}

\begin{proof}
Following the file
\texttt{theories/stable/diag\_bilinear\_tensor.v}, the recipe is pure
structural composition through the tensor--hom adjunction
\autoref{thm:thm512}. Let
$\Phi_{\mathrm{curry}} = \texttt{tensor\_curry}\,\Phi :
\icones_{\mathrm{hom}}\,G\,(E_A \multimap E_B)$. Paper Eq.\ 5.1
(\autoref{lem:eq51}) reads
\[
  \Phi(g \otimes_p a) =
  \texttt{linhom\_fun}\,(\Phi_{\mathrm{curry}}\,g)\,a,
\]
and unfolding the $\mathrm{Ev}$ of $\mathbf{SCones}$
(\autoref{lem:ev}) on a pair recovers the diagonal evaluation as
\[
  g \;\mapsto\; \mathrm{Ev}\bigl(
    \Psi(g),\; K(g)
  \bigr),\qquad
  \Psi(g) \;:=\; \texttt{linhom\_to\_stablehom}\,
                   (\Phi_{\mathrm{curry}}\,g)
\]
in the cartesian-closed packaging
$(B \mathbin{\Rightarrow_s} C) \mathbin{\&} B \to C$. The four
ingredients are then assembled by \autoref{lem:linhom-to-stablehom},
\autoref{lem:meas-stable-comp}, \autoref{lem:spair-meas-stable} and
\autoref{lem:ev}, after rescaling $K$ to operator norm $\leq 1$ to
satisfy the image-ball hypothesis of
\rocq{Icones.stable.compose.meas_stable_comp}, and rescaling back at
the end by linearity of $\mathrm{Ev}$ in its first slot.
\end{proof}

\begin{lemma}[M1 --- the lift $\texttt{linhom\_to\_stablehom}$]%
\label{lem:linhom-to-stablehom}
\rocq{Icones.stable.diag_bilinear_tensor.linhom_to_stablehom}\rocqok
\rocq{Icones.stable.diag_bilinear_tensor.linhom_to_stablehom_meas_stable}\rocqok
\uses{def:linhom-car, def:stablehom, def:scones-hom, thm:scones-ccc}
For integrable cones $B, C$ there is a function
\[
  \texttt{linhom\_to\_stablehom} :
  \texttt{linhom\_car}\,\Ar\,B\,C
  \;\longrightarrow\;
  B \mathbin{\Rightarrow_s} C
\]
sending an operator-norm-$\leq 1$ linear continuous map to its
$0$-extension off the unit ball, packaged as a
$\mathbf{SCones}$ internal-hom element. Viewed as a map between the
\emph{integrable-cone} carriers, the lift is itself measurable stable
(\rocq{Icones.stable.diag_bilinear_tensor.linhom_to_stablehom_meas_stable}):
linearity (\rocq{Icones.stable.diag_bilinear_tensor.linhom_to_stablehom_linear}),
operator-norm bound
(\rocq{Icones.stable.diag_bilinear_tensor.linhom_to_stablehom_bounded}),
$\omega$-continuity on the unit ball
(\rocq{Icones.stable.diag_bilinear_tensor.linhom_to_stablehom_continuous}),
and path preservation
(\rocq{Icones.stable.diag_bilinear_tensor.linhom_to_stablehom_pres_path})
all transfer pointwise from the source $\texttt{linhom\_car}$ through
$\texttt{sc\_clamp}$.
\end{lemma}

\begin{proof}
Pointwise on the unit ball, the lift is the underlying linear function
(\rocq{Icones.stable.diag_bilinear_tensor.linhom_to_stablehom_E});
off the ball it vanishes
(\rocq{Icones.stable.diag_bilinear_tensor.linhom_to_stablehom_off})
by $\texttt{sc\_clamp\_offball}$. Linearity, the $0$-bound, $\omega$-continuity
and path preservation are then verified by a unit-ball case split using
the corresponding pointwise property of $\texttt{linhom\_fun}$. The
measurable-stability assembly is then: linear $\Rightarrow$ totally
monotonic (\autoref{lem:linear-totmono}), linear continuous $+$ the
unit-ball-determined Scott-continuity lift
(\autoref{lem:linear-scott-of-omega}, paper Lemma 2.8/2.10), plus the
path-preservation transfer through $\texttt{sc\_clamp}$.
\end{proof}

\begin{lemma}[M2 --- composition of an $\icones_{\mathrm{hom}}$ with a meas-stable map]%
\label{lem:meas-stable-comp}
\rocq{Icones.stable.diag_bilinear_tensor.meas_stable_comp_post}\rocqok
\rocq{Icones.stable.diag_bilinear_tensor.meas_stable_comp_pre}\rocqok
\uses{lem:ders, thm:730}
An $\icones_{\mathrm{hom}}$ is in particular measurable stable
(\rocq{Icones.homs.icone_cat.icones_meas_stable}, paper Lemma 7.31 at
the morphism level) and norm-decreasing
(\rocq{Icones.icones.icone_cat.cones_hom_norm_le1}). Composing a
measurable stable map with an $\icones_{\mathrm{hom}}$ on either side
lands in the meas-stable closure under composition of
\autoref{thm:730}: the post-composition
$g \circ f$ for $f$ an $\icones_{\mathrm{hom}}$ and $g$ meas-stable is
meas-stable, and the pre-composition $g \circ f$ for $f$ meas-stable
(with the image-ball hypothesis) and $g$ an
$\icones_{\mathrm{hom}}$ is meas-stable.
\end{lemma}

\begin{lemma}[M3 --- the diagonal stable pairing]\label{lem:spair-meas-stable}
\rocq{Icones.stable.diag_bilinear_tensor.id_spair_meas_stable}\rocqok
\rocq{Icones.stable.diag_bilinear_tensor.spair_meas_stable}\rocqok
\uses{lem:meas-stable-comp, thm:scones-prod}
For measurable stable maps $f_1 : G \to A$ and $f_2 : G \to B$, the
cartesian pairing $g \mapsto \texttt{sprod\_pair}\,(f_1\,g)\,(f_2\,g)
: G \to A \mathbin{\&} B$ is measurable stable. The degenerate case
$f_1 = \mathrm{id}_G : G \to G$ paired with a meas-stable $K : G \to A$
is
\rocq{Icones.stable.diag_bilinear_tensor.id_spair_meas_stable}; the
general two-argument pairing
\rocq{Icones.stable.diag_bilinear_tensor.spair_meas_stable} follows by
projecting through $\texttt{scones\_tuple}$. Total monotonicity is the
bool-indexed product (7.1) inequality (\autoref{thm:scones-prod},
\rocq{Icones.stable.scones_cat.cones_prod_le_compI}); boundedness by
the cones product sup; Scott continuity on the unit ball is
componentwise on each $f_i$; path preservation transports through the
product test family.
\end{lemma}

\begin{remark}[The Bang-level instantiation]
With $E_A := \mathord{!}A$ and $E_B := \mathord{!}B$
(\autoref{def:bang-object}, \autoref{ch:exponential}), the bridge
\autoref{thm:diag-bilin} reads: for a meas-stable
$K : G \to \mathord{!}A$ and a bilinear
$\Phi : G \otimes \mathord{!}A \to \mathord{!}B$, the diagonal
$g \mapsto \Phi(g \otimes K(g))$ is meas-stable. This is the structural
unblocker of \emph{generic} CBV recursion at $\mathtt{ne\_fix}$
(\autoref{sec:cbv-ppl}): the body $M$ of
$\mathtt{fix\,s.\,\lambda x.\,\ldots}$ takes the recursive call
$K$ and the captured environment $g$ and assembles them through a
bilinear-into-tensor $\Phi$. The bridge of this section gives the
meas-stability of the diagonal at \emph{every} body $M$, unblocking
the CBV side of the open issue tracked as \texttt{docs/PPL.md \S 7}.
(The earlier CBV mass-one identity sidestepped the bridge through a
manual Kleene chase at this single example; that proof has not yet
been ported to the new linhom-valued \texttt{ppl\_cbv.v}
interpretation.)
\end{remark}

%%%%%%%%%%%%%%%%%%%%%%%%%%%%%%%%%%%%%%%%%%%%%%%%%%%%%%%%%%%%%%%%%%%%%%%%
\section{CBN arithmetic on
\texorpdfstring{$\FMeas$}{FMeas}}\label{sec:cbn-arith}

\paragraph{Why the genuine arithmetic.}
The lightweight option $\gamma$ of \texttt{ppl\_cbn\_eff.v} interprets
$\mathtt{ne\_add}/\mathtt{ne\_mul}$ as the degenerate-at-zero clauses
(the headlines baseline of
\texttt{theories/programs/ppl\_cbn\_headlines.v}); the genuine
arithmetic, packaged at the $\FMeas$ math foundation level, lives in
\texttt{theories/programs/ppl\_cbn\_arith.v}. The relationship between
the two is the same as that between $\Phi$ unrestricted on the tensor
and $\Phi$ on the cartesian pair: only the bridge of
\autoref{sec:diag-bilin} makes the genuine arithmetic stable as a
$\mathbf{SCones}$ map of two arguments.

\begin{definition}[The $\FMeas$ arithmetic foundation]\label{def:cbn-arith}
\rocq{Icones.programs.ppl_cbn_arith.add_FMeas}\rocqok
\rocq{Icones.programs.ppl_cbn_arith.mul_FMeas}\rocqok
\uses{thm:cbv-fmeas-lax, lem:eq51}
The genuine arithmetic operations on the $\FMeas$ cone are
$\mathtt{add\_FMeas},\,\mathtt{mul\_FMeas} :
\FMeas(\Robj) \otimes \FMeas(\Robj) \to \FMeas(\Robj)$, the bilinear
maps obtained by post-composing the FMeas lax-monoidal map
\rocq{Icones.homs.fmeas_lax.fmeas_lax} of
\autoref{thm:cbv-fmeas-lax} with the $\FMeas$-functorial action of the
measurable $\mathbb{R}\times\mathbb{R}\to\mathbb{R}$ arithmetic. Each
is an $\icones_{\mathrm{hom}}$ out of the \emph{tensor} ---
$\Phi : \FMeas(\Robj) \otimes \FMeas(\Robj) \to \FMeas(\Robj)$ ---
and so cannot be installed verbatim as the CBN clause whose source is
the cartesian $\sem{\Gamma}_n$. The bridge
\autoref{thm:diag-bilin} is exactly the missing ingredient: it gives
the meas-stability of $\gamma \mapsto \Phi(\sem{M}_n(\gamma) \otimes
\sem{N}_n(\gamma))$ on the cartesian $\sem{\Gamma}_n$, for any
meas-stable $\sem{M}_n, \sem{N}_n$.
\end{definition}

\begin{remark}[The CBN \texttt{add}/\texttt{mul} install]\label{rem:cbn-arith-install}
Installing the refined CBN $\mathtt{add}/\mathtt{mul}$ clauses on top of
\autoref{def:cbn-arith} is a direct consequence of
\autoref{thm:diag-bilin} (twice, with $K := \sem{N}_n$ and
$\Phi := \mathtt{add\_FMeas}$ / $\mathtt{mul\_FMeas}$), once the
$\gamma$-degenerate clauses of \texttt{ppl\_cbn\_eff.v} are swapped
out. We have not packaged the refined install (the headlines of
\texttt{ppl\_cbn\_headlines.v} are stated under option $\gamma$ and
that is what feeds the dashboard right now), but the math is now in
place and the swap is structural.
\end{remark}

%%%%%%%%%%%%%%%%%%%%%%%%%%%%%%%%%%%%%%%%%%%%%%%%%%%%%%%%%%%%%%%%%%%%%%%%
\section{The CBN headline:
\texorpdfstring{$\mathtt{ex\_geom}$}{ex\_geom} has mass one}%
\label{sec:cbn-geom}

\paragraph{The geometric program, CBN-side.}
The CBN-side mass-one identity for $\mathtt{ex\_geom}$ is the focus
of this section. (The CBV-side parallel is
\rocq{Icones.programs.ex_reject_headline.ex_geom_cbv_mass_one} of
\autoref{thm:cbv-ex-geom}, proved through the seeded value-fixpoint
combinator \rocq{Icones.programs.infra.em_fix_value.fix_comb} of
\autoref{sec:cbv-yfix-fun}.) The surface term is
unchanged:
\[
  \mathtt{ex\_geom}
  \;:=\;
  (\mathtt{fix\,"g":::tfun\,tunit\,tR\,in}\,
   \backslash"\_":::tunit \,\Rightarrow\,
   \mathtt{if\;Bernoulli\{1/2\}\;then\;[|0|]\;else\;[|1|] + \#"g"\,@\,()})\,@\,(),
\]
the geometric counter.

\paragraph{The CBN-side reduction equation, unary.}
A direct reading of the body in CBN gives the recursive-value equation
in $\FMeas(\Robj)$:
\[
  \mu \;\mapsto\;
  \tfrac{1}{2}\cdot \delta_0 \;+\; \tfrac{1}{2}\cdot
  \texttt{shift\_FMeas}\,1\,\mu,
\]
where $\mu$ is the candidate value of the recursive call
$\#"g"\,@\,()$ and $\texttt{shift\_FMeas}\,1$ is the pushforward
through the measurable shift $(+1)$. \textbf{The crucial point}: the
ELSE branch $[|1|] + \#"g"\,@\,()$ is \emph{not} bilinear in two
arguments, but \emph{unary} in $\mu$: the first slot is the constant
Dirac $\delta_1$. So the operation collapses to a unary linear
pushforward through $(+1)$ --- an $\icones_{\mathrm{hom}}$, hence a
$\mathbf{SCones}$ morphism through the dereliction
\autoref{lem:ders} --- and \textbf{no bilinear bridge is required at
the headline}. This is what makes the closure available
\emph{independently} of \autoref{sec:diag-bilin}.

\begin{definition}[The unary shift on $\FMeas$]\label{def:cbn-shift}
\rocq{Icones.programs.ppl_cbn_geom.shift_meas}\rocqok
\rocq{Icones.programs.ppl_cbn_geom.shift_lift}\rocqok
\rocq{Icones.programs.ppl_cbn_geom.shift_scones}\rocqok
\uses{lem:ders}
For $d : \mathbb{R}$, the measurable shift $c \mapsto c + d$ packages
as an $\mathbf{Ar}$ morphism
\rocq{Icones.programs.ppl_cbn_geom.shift_meas},
lifts to an $\icones_{\mathrm{hom}}\,\FMeas(\Robj)\,\FMeas(\Robj)$ via
$\FMeas$-functoriality
(\rocq{Icones.programs.ppl_cbn_geom.shift_lift}), and promotes to a
$\mathbf{SCones}$ morphism through the dereliction
$\mathrm{Der}$ (\autoref{lem:ders})
as \rocq{Icones.programs.ppl_cbn_geom.shift_scones}. The Dirac identity
\rocq{Icones.programs.ppl_cbn_geom.shift_lift_dirac} reduces
$\texttt{shift\_lift}\,d\,\delta_r = \delta_{d+r}$, and the mass identity
\rocq{Icones.programs.ppl_cbn_geom.shift_lift_setT} that
$\texttt{shift\_lift}\,d$ preserves total mass.
\end{definition}

\begin{definition}[The CBN-side geometric body operator]\label{def:phi-cbn-geom}
\rocq{Icones.programs.ppl_cbn_geom.phi_CBN_geom}\rocqok
\rocq{Icones.programs.ppl_cbn_geom.phi_CBN_geom_E}\rocqok
\uses{def:cbn-shift, rem:cbn-bool, def:cbv-bool-case-hom, lem:ev}
The CBN-side operator
$\phi_{\mathrm{CBN}}^{\mathrm{geom}} :
\mathbf{SCones}(\FMeas(\Robj),\FMeas(\Robj))$ realises the body's
reduction equation as a $\mathbf{SCones}$ endomorphism of
$\FMeas(\Robj)$. It is the
$\texttt{spair}$-then-$\mathrm{Ev}$ composite
\[
  \phi_{\mathrm{CBN}}^{\mathrm{geom}} \;=\;
  \mathrm{Ev} \;\circ\; \texttt{spair}\,(
    [\![\,\texttt{bernoulli}\,;\;\delta_0\,;\;\texttt{shift\_lift}\,1\,]\!]
  )\,\mathrm{id},
\]
where the $\texttt{linhom\_car}$ co-pairing
\rocq{Icones.programs.ppl_cbn_geom.phi_geom_if_linhom} is the
\rocq{Icones.programs.infra.bool_case_hom.bool_case_linhom} of
\autoref{def:cbv-bool-case-hom} applied to the two branches
(constant $\delta_0$, unary shift $\texttt{shift\_scones}\,1$),
$\texttt{ders}$-promoted to a $\mathbf{SCones}$ morphism out of
$\mathtt{bool\_cone\_car}$, pre-composed with the constant
$\texttt{Bernoulli}\,(1/2)$ scrutinee, and then diagonal-evaluated by
\autoref{lem:ev}. The pointwise reduction on the unit ball is
\rocq{Icones.programs.ppl_cbn_geom.phi_CBN_geom_E}:
\[
  \texttt{sc\_fun}\,\phi_{\mathrm{CBN}}^{\mathrm{geom}}\,\mu
  \;=\;
  \texttt{bool\_case}\,(\mathtt{Bernoulli}\,1/2)\,
    (\delta_0)\,(\texttt{shift\_lift}\,1\,\mu),
\]
i.e.\ the equation $\mu \mapsto \tfrac{1}{2}\delta_0 +
\tfrac{1}{2}\texttt{shift\_lift}\,1\,\mu$.
\end{definition}

\paragraph{The Kleene cascade closed form.}
The mass identity now reduces to the Kleene cascade for
$\phi_{\mathrm{CBN}}^{\mathrm{geom}}$ from $\bot$, using the same
per-iterate cascade pattern that the (now-retired) CBV side used
through \texttt{em\_fix\_arr.v}. The per-iterate mass is closed in
three steps.

\begin{lemma}[Per-iterate cascade]\label{lem:cbn-kleene-cascade}
\rocq{Icones.programs.ppl_cbn_geom.kleene_geom_S_E}\rocqok
\rocq{Icones.programs.ppl_cbn_geom.kleene_geom_S_mass}\rocqok
\rocq{Icones.programs.ppl_cbn_geom.kleene_geom_mass_closed}\rocqok
\uses{def:phi-cbn-geom}
Define $\texttt{kleene\_geom}\,n := \texttt{kleene}\,(\texttt{sc\_fun}\,
\phi_{\mathrm{CBN}}^{\mathrm{geom}})\,n : \FMeas(\Robj)$ (the
$\bot$-seeded $n$-th Kleene iterate). On the unit ball,
\[
  \texttt{kleene\_geom}\,(n+1) =
  \tfrac{1}{2}\cdot \delta_0 \;+\;
  \tfrac{1}{2}\cdot \texttt{shift\_lift}\,1\,(\texttt{kleene\_geom}\,n)
\]
(\rocq{Icones.programs.ppl_cbn_geom.kleene_geom_S_E}), whence the mass
recurrence
$\texttt{mass}(\texttt{kleene\_geom}\,(n+1)) =
\tfrac{1}{2} + \tfrac{1}{2}\cdot
\texttt{mass}(\texttt{kleene\_geom}\,n)$
(\rocq{Icones.programs.ppl_cbn_geom.kleene_geom_S_mass}) and the
closed-form
\[
  \texttt{mass}(\texttt{kleene\_geom}\,n) = 1 - (1/2)^n
\]
(\rocq{Icones.programs.ppl_cbn_geom.kleene_geom_mass_closed}).
\end{lemma}

\begin{proof}
The first equation rewrites
$\texttt{kleene\_geom}\,(n+1) = \texttt{sc\_fun}\,
\phi_{\mathrm{CBN}}^{\mathrm{geom}}\,(\texttt{kleene\_geom}\,n)$ via
\rocq{Icones.programs.ppl_cbn_geom.kleene_geom_S} and unfolds
$\texttt{sc\_fun}\,\phi_{\mathrm{CBN}}^{\mathrm{geom}}$ on the unit
ball via \rocq{Icones.programs.ppl_cbn_geom.phi_CBN_geom_E}, applicable
because \rocq{Icones.programs.ppl_cbn_geom.kleene_geom_ball} maintains
the unit-ball invariant. The mass recurrence is bilinearity of
$\FMeas$-mass plus the Dirac mass $1$
(\rocq{Icones.mcones.fmeas.dirac_fmeas_setT_E}) and the shift's
mass-preservation \rocq{Icones.programs.ppl_cbn_geom.shift_lift_setT}.
The closed-form is induction on $n$.
\end{proof}

\begin{lemma}[Mass convergence]\label{lem:cbn-kleene-cvg}
\rocq{Icones.programs.ppl_cbn_geom.kleene_geom_mass_cvg}\rocqok
\uses{lem:cbn-kleene-cascade}
$\texttt{mass}(\texttt{kleene\_geom}\,n) \to 1$ as $n \to \infty$
(in $\overline{\mathbb{R}}$), since $(1/2)^n \to 0$.
\end{lemma}

\begin{theorem}[CBN headline: $\mathtt{ex\_geom}$ has mass one]%
\label{thm:cbn-ex-geom}
\rocq{Icones.programs.ppl_cbn_geom.ex_geom_CBN_fix}\rocqok
\rocq{Icones.programs.ppl_cbn_geom.ex_geom_CBN_mass_one}\rocqok
\uses{def:phi-cbn-geom, lem:cbn-kleene-cvg, lem:cbn-recursion}
The least fixpoint
$\texttt{ex\_geom\_CBN\_fix} := \texttt{sfix}\,
\phi_{\mathrm{CBN}}^{\mathrm{geom}} : \FMeas(\Robj)$
of the geometric body operator has total mass $1$ in
$\overline{\mathbb{R}}$:
\[
  \texttt{fmeas\_mu}\,(\texttt{ex\_geom\_CBN\_fix})\,[\,\mathit{setT}\,]
  \;=\;
  1.
\]
\end{theorem}

\begin{proof}
By definition of the $\mathbf{SCones}$ least fixpoint
\rocq{Icones.stable.fixpoint.sfix}, $\texttt{ex\_geom\_CBN\_fix}$
unfolds to the $\texttt{cone\_sup\_ball}$ of the Kleene chain
$\texttt{kleene\_geom}$. The HB instance identifies this with the
$\FMeas$-side $\texttt{fmeas\_sup\_ball}$, whose mass on the full set
agrees with the limit of the iterate masses
(\rocq{Icones.mcones.fmeas.fmeas_sup_ballE} and
\rocq{Icones.mcones.fmeas.fmeas_sup_cvg}). The limit is $1$ by
\autoref{lem:cbn-kleene-cvg} and the uniqueness of limits in the
Hausdorff space $\overline{\mathbb{R}}$.
\end{proof}

\paragraph{The CBN headline avoids the bridge.}
The key insight that makes the CBN headline closable in $\sim 700$
lines, independently of \autoref{sec:diag-bilin}, is the unary
``shift-by-Dirac'' observation of the prologue: at the headline
$\mathtt{ex\_geom}$, the body's ELSE branch is unary in the recursive
value, the first slot of the would-be $\mathtt{add}$ being a constant
Dirac. The diagonal bilinear bridge of \autoref{sec:diag-bilin}
remains the structural unblocker for \emph{generic} bilinear
arithmetic on the recursive value --- the refined CBN
$\mathtt{add}/\mathtt{mul}$ install of
\autoref{rem:cbn-arith-install} above --- and the structural unblocker
of generic CBV recursion at $\mathtt{ne\_fix}$
(\autoref{sec:cbv-ppl}).

%%%%%%%%%%%%%%%%%%%%%%%%%%%%%%%%%%%%%%%%%%%%%%%%%%%%%%%%%%%%%%%%%%%%%%%%
\section{Status}\label{sec:cbn-status}

\paragraph{Delivered and axiom-free.}
Everything in this chapter is axiom-clean (only the three classical
\texttt{boolp} axioms): the CBN type / context translations
$\mathtt{tyD\_CBN}$ / $\mathtt{ctxD\_CBN}$
(\rocq{Icones.programs.ppl_cbn.tyD_CBN},
\rocq{Icones.programs.ppl_cbn.ctxD_CBN},
\autoref{def:cbn-tyD}); the pure fragment of the term denotation
\rocq{Icones.programs.ppl_cbn.eD_CBN} with the fixpoint reduction
\rocq{Icones.programs.ppl_cbn.eD_CBN_fix_E}
(\autoref{def:cbn-eD}, \autoref{lem:cbn-recursion}); the CBN
boolean cascade through the $\texttt{sc\_to\_sh}$ promotion
(\autoref{rem:cbn-bool},
\rocq{Icones.programs.infra.bool_case_scones.bool_case_scones},
\rocq{Icones.programs.ppl_cbn_bool.cbn_if_clause_def}); the CBN-side
$\FMeas$ arithmetic foundation
\rocq{Icones.programs.ppl_cbn_arith.add_FMeas} /
\rocq{Icones.programs.ppl_cbn_arith.mul_FMeas} (\autoref{def:cbn-arith});
the diagonal bilinear stability bridge
\rocq{Icones.stable.diag_bilinear_tensor.meas_stable_diag_bilinear_tensor}
with the four milestone lemmas
\rocq{Icones.stable.diag_bilinear_tensor.linhom_to_stablehom_meas_stable}
/
\rocq{Icones.stable.diag_bilinear_tensor.meas_stable_comp_post} /
\rocq{Icones.stable.diag_bilinear_tensor.id_spair_meas_stable} /
\rocq{Icones.stable.diag_bilinear_tensor.spair_meas_stable}
(\autoref{thm:diag-bilin} and \autoref{lem:linhom-to-stablehom} /
\autoref{lem:meas-stable-comp} / \autoref{lem:spair-meas-stable});
the CBN headline
\rocq{Icones.programs.ppl_cbn_geom.ex_geom_CBN_mass_one}
(\autoref{thm:cbn-ex-geom}) and its Kleene cascade
\rocq{Icones.programs.ppl_cbn_geom.kleene_geom_mass_closed} /
\rocq{Icones.programs.ppl_cbn_geom.kleene_geom_mass_cvg}.

\paragraph{Honest scope notes.}
The CBV / CBN soundness theorem (a proof that
$\sem{-}_{\mathrm{CBV}}$ and $\sem{-}_n$ agree on a suitable observable)
is \emph{not} delivered here; bridging the two readings at the level
of the $\mathord{!}$-coalgebra would require commuting $\mathord{!}$
with the effect-bearing types, which has no closed-form realisation in
the SAFT-built $\mathord{!}$. The refined CBN $\mathtt{add}/\mathtt{mul}$
install on top of \autoref{def:cbn-arith}, via
\autoref{thm:diag-bilin}, is a structural consequence of what is
in place; we have not packaged it because the headlines of
\texttt{theories/programs/ppl\_cbn\_headlines.v} are stated under
option $\gamma$ (lightweight degenerate arithmetic), which is what
feeds the present mass identities. The genuine score in CBN under the
\emph{option-B} translation $\sem{\mathtt{tunit}}_n = \top$ collapses
by terminal uniqueness; a faithful CBN score requires a different
translation of $\mathtt{tunit}$ (e.g.\ $\sem{\mathtt{tunit}}_n =
\FMeas(\ast)$), the chase of which is orthogonal to this chapter.
Mutual recursion at free-coalgebra types is out of scope.

\paragraph{The two centrepieces and the relationship to Chapter
\ref{ch:cbv}.}
The two structural results of this chapter ---
\autoref{thm:diag-bilin} (the bridge) and
\autoref{thm:cbn-ex-geom} (the CBN headline) --- form a coherent pair.
The CBN headline establishes that the geometric program has mass $1$
under the $\mathbf{SCones}$ / co-Kleisli reading; the bridge establishes
the structural unblocker that, when packaged, would deliver the
identity \emph{generically} (not example-by-example), for arbitrary
bilinear arithmetic and recursive bodies. The headline is an
example-specific shortcut that the bridge would generalise; the
bridge is the math content that the headline avoids by exploiting the
geometric program's specific shape (Bernoulli + Dirac constants). The
CBV-side mass-one identity has landed:
\rocq{Icones.programs.ex_reject_headline.ex_geom_cbv_mass_one}
(\autoref{thm:cbv-ex-geom}), proved on the seeded value-fixpoint
combinator of \autoref{sec:cbv-yfix-fun}, so the two
interpretations' identities now form a coherent triple with the
bridge: the CBN headline (\autoref{thm:cbn-ex-geom}), the CBV twin
(\autoref{thm:cbv-ex-geom}), and the structural bridge
(\autoref{thm:diag-bilin}) that relates the bilinear arithmetic of
the two readings. All formalisation
content lives in the axiom-free core.
